\begin{table}[!htbp]
\centering
\caption{Fund definitions on the broad-measure event sample}\label{tab:holdings_coverage}
\small
\setlength{\tabcolsep}{5pt}
\begin{tabular}{llrrrrr}
\toprule
Type & Period & Events & Strict & + variants & Broad & Mixed count \\
\midrule
Addition & 2007--11 & 8 & 0.10 & 0.15 & 0.15 & 0.0 \\
Addition & 2012--16 & 27 & 0.37 & 0.49 & 0.52 & 0.4 \\
Addition & 2017--21 & 58 & 0.92 & 1.02 & 1.11 & 1.0 \\
Addition & 2022--26 & 78 & 0.41 & 0.78 & 0.97 & 8.8 \\
Deletion & 2007--11 & 7 & 0.20 & 0.26 & 0.26 & 0.0 \\
Deletion & 2012--16 & 25 & 0.31 & 0.46 & 0.48 & 0.1 \\
Deletion & 2017--21 & 40 & 0.83 & 0.92 & 0.99 & 1.0 \\
Deletion & 2022--26 & 68 & 0.52 & 0.93 & 1.18 & 9.4 \\
\bottomrule
\end{tabular}
\begin{minipage}{\linewidth}\footnotesize\vspace{5pt}
Demand is percent of shares outstanding, sign-adjusted for deletions. All three demand definitions in a row use the same retained event sample. Strict comprises exact benchmark followers; the next column adds related benchmark variants; Broad also includes the benchmark-linked mixed-fund category. Mixed count is the mean number of reporting mixed funds per event, not their ownership percentage. These data do not measure the complete investor universe.
\end{minipage}
\end{table}
